%% ma: L10-B14-0101
%% lop: 10
%% chuong: 5
%% bai: 14
%% dang: 10.14.8
%% loai: TLN
%% muc: VD
%% dap_an: "1,87"
%% nguon: "chuyen-de-cac-so-dac-trung-cua-mau-so-lieu-khong-ghep-nhom-toan-10-chuong-trinh-moi.pdf (GV. Phan Nhật Linh) (P05), Bài 03 Dạng 2 Phần III Câu 7"
%% kien_thuc: "Lớp 10 Bài 13 (số trung bình), Bài 14 (phương sai, độ lệch chuẩn)"
%% trang_thai: chinh-thuc
%% ghi_chu: "Đáp số gốc 1,874 có 5 ký tự, đã yêu cầu làm tròn đến hàng phần trăm. Tự duyệt theo yêu cầu giáo viên 22/10/2026."
\begin{ex}
Cho bảng điểm kiểm tra môn Toán cuối học kì 2 (thang điểm 10) của $40$ học sinh lớp 10C như sau:
\begin{center}
\begin{tabular}{|c|c|c|c|c|c|c|c|}\hline
Điểm & $5$ & $6$ & $7$ & $8$ & $9$ & $10$ & Cộng\\\hline
Tần số & $5$ & $12$ & $8$ & $9$ & $4$ & $2$ & $40$\\\hline
\end{tabular}
\end{center}
Tính phương sai $s^2$ của mẫu số liệu (làm tròn đến hàng phần trăm).
\shortans{1,87}
\loigiai{$\overline x=\dfrac{281}{40}=7{,}025$; $s^2=\dfrac1{40}\sum n_i\left(x_i-7{,}025\right)^2\approx1{,}874\approx1{,}87$.}
\end{ex}
